\chapter*{Notation}
\addcontentsline{toc}{chapter}{Notation}  
%
%
The natural unit system $\hbar = c = k_{B} = 1$, and thus $h = 2\pi$, is used throughout the thesis.
This implies the dimensional relation [Energy] $=$ [Mass] $=$ [Temperature] $=$ [Time$^{-1}$] $=$ [Length$^{-1}$].
The Fermi coupling constant is $G_F = 1.1663787(6) \times 10^{-5}\,\mathrm{GeV}^{-2}$~\cite{pdg}.
The vacuum expectation value (vev) of the Standard Model Higgs is taken to be $v = 174.104\,\mathrm{GeV}$.
The metric with signature $\text{diag}(-1,+1,+1,+1)$ is used.
The reduced Planck mass is $M_{\mathrm{Planck}}= \sqrt{\hbar c/(8\pi G_{\text{Newton}})} = \sqrt{1/(8\pi G_{\text{Newton}})} = 2.435 \times 10^{18}\,\myunit{GeV}$.
A variable written in bold format, e.g. $\mathbf{x}$, represents a tuple or vector.
Implicit summation over repeated indices is assumed.








%%%%%%%%%%%%%%%%%%%%%%%%%%%%%%%%%%%%%%% END %%%%%%%%%%%%%%%%%%%%%%%%%%%%%%%%%%%%%%%%%%%%%%%%%%%%